\section{Performance analysis}\label{sec:performance}

\subsection{Convergence and efficiency of the integrators}\label{sec:performance-convergence}

\begin{figure}[tb]
  \centering
  \maybefig{fig_convergence}{\linewidth}
  \caption{Integration back-ends at a near-boundary unstable point of the
  showcase system, where the exact count $\Zint=2$ is known, so the integer
  residual $\varepsilon$ is an \emph{absolute} error measure requiring no
  reference. (a)~Accuracy versus requested tolerance; the dotted line marks
  the practical rule $\varepsilon\approx100\,\mathrm{tol}$. (b)~Cost versus
  tolerance. (c)~Efficiency front. Error is on the vertical axis in both (a)
  and (c).}\label{fig:convergence}
\end{figure}

\begin{table}[tb]
  \centering
  \caption{Integer residual (left block) and single-point CPU time (right
  block) of the back-ends at selected tolerances; system and point as in
  Fig.~\ref{fig:convergence}.}\label{tab:convergence}
  \maybeinput{tables/tab_convergence.tex}
\end{table}

Figure~\ref{fig:convergence} and Table~\ref{tab:convergence} quantify the
choice of back-end, and three points deserve emphasis. First, all adaptive
back-ends track the requested tolerance over the studied range, with the
integer residual settling at roughly two orders of magnitude above it -- the
rule of thumb $\varepsilon\approx100\,\mathrm{tol}$ quoted in
Section~\ref{sec:integerness}. Second, the order of the pair matters far
more than its family. The classical fifth-order pair is perfectly adequate
at the default $10^{-5}$, but at $10^{-3}$ it under-resolves the
$\omega^{-1}$ ripple of this velocity-feedback system so badly that
$\Zraw$ is off by $0.8$ -- within a hair of losing the count -- whereas the
ninth-order default is three orders of magnitude more accurate there and,
at tolerances below $10^{-6}$, also cheaper. The linearly implicit
Rosenbrock pair is uniformly dominated, as expected once one recognizes that
the problem has no stiffness for it to exploit
(Section~\ref{sec:phaseode}). Third, at the practically relevant setting
$10^{-5}$ the back-ends are within a small factor of one another in both
accuracy and cost: \emph{the choice of integrator is simply not critical
there}, and the default is chosen for graceful degradation rather than for
speed.

\begin{table}[tb]
  \centering
  \caption{Approaching a Hopf boundary of the showcase system at distance
  $\Delta P$ in the gain. The peak half-width equals $|\sest|$; the marching
  solver silently loses one root once the peak falls below its step
  resolution, while the integer residual stays small
  ($\Zraw$ shifts by exactly one).}\label{tab:boundarystress}
  \maybeinput{tables/tab_boundary_stress.tex}
\end{table}

The one situation in which the back-ends genuinely differ is the immediate
vicinity of a stability boundary, and the outcome is the opposite of what
the folklore about adaptive solvers suggests
(Table~\ref{tab:boundarystress}). As the crossing root approaches the
integration line, the integrand peak narrows in proportion to $|\sest|$; at
$\Delta P=10^{-8}$ it is $6.6\cdot10^{-10}$ wide and $1.5\cdot10^{9}$ tall,
while $10^{-3}$ away from it the integrand is back to its smooth
background. The marching solver then steps clean over the peak, loses
exactly $\pm\pi$, reports one root too few -- and returns a success flag and
a perfectly healthy integer residual. Gauss--Kronrod survives the same test,
though not by cleverness: the persistent delay ripple forces it to subdivide
finely across the whole range anyway, and that global refinement happens to
resolve the spike. Neither method carries a guarantee, since both sample
pointwise. What does carry information is the tracked root: at
$\Delta P=10^{-8}$ it still reports $\sest=+6.6\cdot10^{-10}$, correctly
placing a root to the right of the axis, so the sign test of
Section~\ref{sec:integerness} exposes the inconsistency for free. This is
the practical recommendation: use quadrature when only $\Zint$ is wanted,
use the ordered march when the root is wanted too -- and let the two outputs
of the march check each other.

\subsection{Chart-level cost, memory and cold start}\label{sec:performance-cost}

\begin{table}[tb]
  \centering
  \caption{Warm chart-level costs (median of repeated runs,
  \NumThreads{}~threads): full brute-force sweeps and MDBM boundary tracing
  for three systems of increasing evaluation cost, with allocated
  memory.}\label{tab:gridtimings}
  \maybeinput{tables/tab_grid_timings.tex}
\end{table}

Table~\ref{tab:gridtimings} reports the cost of a \emph{complete}
$100\times100$ chart -- ten thousand parameter points -- on a
\CPUmodel{} (\NumThreads{} threads, \RAMgb\,GB RAM, Julia~\JuliaVersion,
package revision \PkgCommit), with the per-point cost listed alongside so
that other resolutions can be scaled directly. For scalar characteristic
functions a full chart takes a few tens of milliseconds, and the entire
hybrid workflow (colored sweep plus MDBM boundary) finishes well inside one
second -- genuinely interactive. For the showcase DDAE every evaluation runs
the automatic extraction of Section~\ref{sec:extraction}, i.e.\ a $7\times7$
dual-number determinant, and the chart takes on the order of a minute at
this resolution; the convenience of never deriving $\Dfun$ by hand thus
costs roughly an order of magnitude against a hand-coded characteristic
function, which the same table quantifies. Memory stays modest throughout
because the phase ODE carries a single scalar state and the sweep stores
only the five result fields; no large matrices are formed at any point.
Cold start is a one-time cost dominated by package loading: in a fresh
Julia process, loading takes about $5.5$\,s and the first thousand-point
sweep $0.1$--$0.3$\,s thanks to the precompilation workload shipped with the
package, after which every further chart runs at the warm speeds of
Table~\ref{tab:gridtimings}.

\begin{table}[tb]
  \centering
  \caption{What the MDBM boundary trace buys: from an initial mesh with the
  listed refinement levels, the traced boundary has the equivalent resolution
  of a uniform grid of the given size, at the listed number of function
  evaluations. The saving factor is the ratio of the two point
  counts.}\label{tab:mdbm}
  \maybeinput{tables/tab_mdbm.tex}
\end{table}

Table~\ref{tab:mdbm} makes the boundary-tracing argument of
Section~\ref{sec:mdbm} concrete. Four refinement levels from a
$20\times20$ mesh yield a boundary resolved as finely as a
$320\times320$ uniform grid would resolve it, while evaluating only the
cells that actually bracket the boundary -- a saving of two orders of
magnitude in function evaluations, which is precisely why the high-resolution
boundary is affordable on top of a deliberately coarse colored background. These numbers substantiate the practical claim
of the paper: the stability chart of a nontrivial delayed system refreshes
fast enough to be explored \emph{interactively}, which changes how such
charts are used in design -- parameter studies and optimization loops can
call the entire chart pipeline as an inner function.

\subsection{Comparison with semi-discretization}\label{sec:performance-sd}

\begin{table}[tb]
  \centering
  \caption{Comparison with semi-discretization
  (\pkg{SemiDiscretizationMethod.jl}, orders 1 and 2, $n$ steps per delay) on
  the reduced showcase system: error of the dominant characteristic exponent
  and CPU time per parameter point. The reference is the dominant root
  refined to machine precision; the two independent formulations agree to
  seven digits.}\label{tab:semidisc}
  \maybeinput{tables/tab_semidisc.tex}
\end{table}

Semi-discretization~\cite{insperger2011semidiscretization} -- here in its
efficient library implementation with the multiplication-free
variant~\cite{bachrathy2026semidisc,bachrathy2026multiplication} --
approximates the dominant characteristic exponent through the spectral
radius of a finite transition matrix, converging polynomially in the number
of discretization steps per delay. Table~\ref{tab:semidisc} compares it with
the present method on the reduced showcase system at a stable, a
near-boundary and an unstable point. The dominant root of the present method
is obtained by the multi-root recipe of Section~\ref{sec:tracking}: ten
tracked minima, all refined, maximal real part. That two entirely
independent formulations agree to seven digits is the most valuable outcome
of the comparison -- it validates both.

Beyond that agreement the two methods are not commensurable, and it would be
misleading to plot them on shared error axes: their accuracy is governed by
different dials. This is the substantive difference. Semi-discretization
requires the user to choose the number of steps per delay \emph{by hand};
the discretization is not adaptive, and an insufficient resolution returns a
confidently wrong exponent with no warning and no internal indication that
anything is amiss. The present method sets its own frequency steps from the
behaviour of the system, exposes the accuracy of the root through the
requested tolerance (Section~\ref{sec:integerness}), and ships two
independent consistency checks: the integer residual, which certifies the
quadrature, and the sign agreement between $\sest$ and $\Zint$, which
certifies the count. A user who wants a more accurate root turns one dial
and can verify the result; a user of the discretization approach must
re-run at several resolutions and compare. The comparison is also
intentionally asymmetric in what the methods deliver: semi-discretization
returns the full dominant multiplier cluster and extends unchanged to
time-periodic systems, where the argument-principle formulation does not
apply (Section~\ref{sec:conclusion}); the present method returns exactly
what an \emph{autonomous} stability chart needs -- the root count, the
boundary, and the dominant root on demand -- at interactive speed.

\subsection{A broader solver comparison}\label{sec:performance-zoo}

\begin{figure}[tb]
  \centering
  \maybefig{fig_solver_zoo}{0.7\linewidth}
  \caption{Error--cost Pareto fronts of six integration back-ends on a
  $40\times40$ sweep of the fourth-order benchmark oscillator
  (\ref{app:gallery}): fixed-step trapezoid (step-count sweep), adaptive
  Gauss--Kronrod, and four ODE solvers (tolerance
  sweeps).}\label{fig:zoo}
\end{figure}

Figure~\ref{fig:zoo} extends the comparison to a wider selection of ODE
solvers on the fourth-order benchmark system. The qualitative picture is
robust: (i)~the fixed-step trapezoid is unbeatable for crude, high-speed
scans but its accuracy stagnates unless the step count grows
with the requested accuracy; (ii)~all adaptive methods follow their tolerance; (iii)~the
lightweight explicit pairs (BS3, Tsit5) are the fastest of the ODE solvers
at moderate accuracy,
while the stiff and auto-switching solvers pay a small overhead that buys
dependable peak capture; the auto-switching default inherits the better of
both regimes. The spread between the methods is far smaller than the gap
that separates all of them from methods that rebuild discretization
matrices at every parameter point, which is the practical message of this
section.